% sections/related_work.tex
% Note: this paper uses \sfrac{a}{b} from xfrac, NOT \nicefrac{a}{b} (nicefrac.sty unavailable).
\section{Related work}
\label{sec:related}

\paragraph{Heavy-tailed policies in RL.}
Tokutake and Yamashita~\cite{tokutake2019} introduced a Student-$t$
actor-critic for robot control. Garg
et~al.~\cite{garg2021heavytails} characterised PPO gradients as
heavy-tailed and proposed robust estimators rather than heavy-tailed
policies. To our knowledge, no recent work benchmarks a Student-$t$
policy with state-dependent degrees of freedom as a drop-in
replacement for the Tanh-Gaussian on standard MuJoCo locomotion.

\paragraph{Safe RL and margin-based reward.}
Constrained Policy Optimization~\cite{achiam2017cpo} treats per-step
margins as constraints in a Lagrangian formulation, retaining a
separate task return. We do not use a Lagrangian: the soft-min
margin \emph{is} the objective, not a penalty added to one. The
Hamilton--Jacobi reachability literature treats margins as value
functions to be solved offline; we treat them as in-the-loop
training signals. The closest precursor is the viability-theoretic
formulation of Aubin~\cite{aubin1991}, which we restrict to a single
trajectory-level scalar.

\paragraph{Implementation matters and minimalist variants.}
Engstrom et~al.~\cite{engstrom2020} and Andrychowicz
et~al.~\cite{andrychowicz2021} demonstrated that PPO's algorithmic
lead is largely an artifact of code-level details, but neither
proposes a simpler algorithm: both validate that the existing tricks
contribute. Subsequent work in language-model fine-tuning has shown
pure score-function methods can match clipped surrogates in
discrete-text regimes; we extend this question to continuous-control
locomotion.

\paragraph{Spectral and Frobenius regularisation.}
Yoshida and Miyato~\cite{yoshida2017spectral} introduced spectral
norm regularisation for generalisation; Miyato
et~al.~\cite{miyato2018spectral} popularised it in GAN training.
Frobenius clamping is a coarser, cheaper variant; we use it because
it admits a clean closed-form projection step and yields the
gradient bound of Lemma~\ref{lem:bounded-score}.

\paragraph{What we are not.}
This paper does \emph{not} propose a new exploration mechanism, a
new representation learning approach, or a new sim-to-real
procedure. We intentionally hold all those axes fixed and ask the
narrow question of whether the algorithmic core of policy-gradient
locomotion can be shrunk to a minimal triple. Pretrained visual
encoders such as DINOv2~\cite{oquab2024dinov2} and quadruped
locomotion stacks~\cite{hwangbo2019anymal} are complementary; we
neither use nor compete with them in this paper.
